\chapter{Introduction}

Ce rapport presente une etude de reproduction et d'analyse critique de
l'article \emph{Dial-In LLM: Human-Aligned LLM-in-the-loop Intent Clustering
for Customer Service Dialogues}~\cite{hong2025dialinllm}. Le travail mene
s'inscrit dans le cadre d'un projet tutore en apprentissage automatique et en
science des donnees, dont l'objet est d'examiner en profondeur une methode
recente, d'en evaluer la reproductibilite a partir des ressources publiques
disponibles et d'en proposer une extension experimentale argumentee.

\section{Contexte scientifique}

L'article etudie un probleme de \emph{clustering d'intentions} dans des
dialogues de service client. Cette tache consiste a regrouper automatiquement
des enonces selon l'intention qu'ils expriment. Elle presente une difficulte
particuliere dans le cadre conversationnel, car la proximite lexicale entre
deux phrases ne garantit ni une equivalence fonctionnelle, ni une separation
nette entre intents. Une methode reposant exclusivement sur la proximite
geometrique d'embeddings peut donc produire des partitions insuffisamment
coherentes sur le plan semantique.

L'originalite de \emph{Dial-In LLM} reside dans l'integration explicite d'un
grand modele de langage au sein de la boucle de clustering. Le LLM n'y est pas
mobilise comme simple outil de pretraitement, mais comme composant de decision,
charge d'evaluer la coherence de clusters intermediaires et de generer des noms
d'intentions. Cette architecture rend la methode particulierement interessante
du point de vue de la reproduction, car elle articule plusieurs composantes
dependantes :
\begin{itemize}[leftmargin=*]
  \item un module d'encodage des enonces ;
  \item une boucle iterative de clustering avec recherche de \texttt{K} ;
  \item une strategie d'echantillonnage des clusters ;
  \item un jugement de coherence par LLM ;
  \item un mecanisme de nommage et de fusion finale.
\end{itemize}

La difficulte du travail de reproduction ne tient donc pas uniquement a la
reexecution d'un code existant, mais a la reconstitution coherente de
l'ensemble d'un protocole methodologique.

\section{Objectifs}

Le projet poursuit quatre objectifs complementaires :
\begin{enumerate}[leftmargin=*]
  \item situer la publication dans son etat de l'art ;
  \item reproduire, dans la mesure du possible, les resultats rapportes par les
  auteurs ;
  \item proposer des evaluations complementaires permettant de tester la
  robustesse de la methode ;
  \item concevoir et evaluer des modifications algorithmiques ou
  computationnelles susceptibles d'ameliorer le pipeline.
\end{enumerate}

Au fil du travail, cette demarche a naturellement conduit a structurer
l'analyse autour de trois niveaux :
\begin{itemize}[leftmargin=*]
  \item les resultats rapportes dans le papier ;
  \item une baseline de reproduction construite dans notre environnement ;
  \item une configuration amelioree issue de nos propres modifications.
\end{itemize}

Cette structuration permet de separer clairement ce qui releve de la
reproduction, ce qui releve de l'analyse critique et ce qui constitue une
contribution experimentale propre au projet.

\section{Demarche generale}

Le travail a ete organise en trois etapes principales.

La premiere a consiste a etablir une baseline de reproduction stable. Cette
phase a implique le nettoyage du depot, l'alignement sur les bons jeux de
donnees, la verification des tailles de benchmarks, la stabilisation de
l'execution sur le serveur du projet et l'identification d'une approximation
reproductible du protocole des auteurs.

La deuxieme a consiste a mener une serie d'evaluations complementaires et
d'ablations ciblees. Plutot que de modifier simultanement l'ensemble du
pipeline, plusieurs leviers ont ete etudies separement : le bloc
\texttt{coherence / naming}, le merge final et la politique de selection de
\texttt{K}. Cette decomposition permet d'attribuer les variations de
performance a des composantes techniques identifiables.

La troisieme a consiste a consolider ces resultats dans une campagne
comparative finale, afin de produire une lecture synthetique et scientifiquement
argumentee de l'etat du projet.

\section{Organisation du rapport}

Le document est organise comme suit :
\begin{itemize}[leftmargin=*]
  \item le chapitre~2 presente l'etat de l'art et situe l'article par rapport
  aux travaux les plus proches ;
  \item le chapitre~3 decrit la reproduction du papier, le protocole retenu,
  l'environnement experimental et les resultats obtenus ;
  \item le chapitre~4 rassemble les evaluations complementaires et les
  modifications proposees ;
  \item le chapitre~5 presente la synthese comparative entre le papier, la
  baseline de reproduction et la configuration amelioree ;
  \item le chapitre~6 propose le bilan general du travail.
\end{itemize}

L'objectif du rapport n'est donc pas seulement de presenter des scores, mais de
fournir une analyse rigoureuse des conditions de reproductibilite du papier,
des limites de l'approche et de la portee des ameliorations introduites.
